\pdfoutput=1
\def\PaperVersion{0.1.3-draft}

\documentclass[11pt]{article}
\usepackage[margin=1in]{geometry}
\usepackage{amsmath,amssymb,amsthm}
\usepackage{enumitem}
\usepackage{hyperref}
\hypersetup{hidelinks,
  pdftitle={An explicit initial Theta-datum with mixed reduction over a split prime},
  pdfauthor={Toshihisa Urahigashi}}

\theoremstyle{plain}
\newtheorem{theorem}{Theorem}[section]
\newtheorem{proposition}[theorem]{Proposition}
\newtheorem{lemma}[theorem]{Lemma}
\newtheorem{corollary}[theorem]{Corollary}
\theoremstyle{definition}
\newtheorem{definition}[theorem]{Definition}
\newtheorem{remark}[theorem]{Remark}
\newtheorem{observation}[theorem]{Observation}

\newcommand{\Fmod}{F_{\mathrm{mod}}}
\newcommand{\Vmod}{\mathbb{V}_{\mathrm{mod}}}
\newcommand{\Vbad}{\mathbb{V}^{\mathrm{bad}}}
\newcommand{\Vgood}{\mathbb{V}^{\mathrm{good}}}
\newcommand{\VV}{\mathbb{V}}
\newcommand{\Ind}[1]{\textup{(Ind#1)}}
\newcommand{\SA}{\textup{(SA)}}
\newcommand{\ord}{\operatorname{ord}}
\newcommand{\Norm}{\operatorname{N}}
\newcommand{\Aut}{\operatorname{Aut}}

\title{An explicit initial $\Theta$-datum with mixed reduction over a split prime,
and the local estimates of inter-universal Teichm\"uller theory IV}
\author{Toshihisa Urahigashi}
\date{Version \PaperVersion\ --- \today\\[3pt]
\small DOI \href{https://doi.org/10.5281/zenodo.22537342}{\texttt{10.5281/zenodo.22537342}} (all versions)\qquad ORCID \href{https://orcid.org/0009-0004-0460-6242}{\texttt{0009-0004-0460-6242}}}

\begin{document}
\maketitle

\begin{abstract}
We exhibit an explicit initial $\Theta$-datum in the sense of
\cite[Definition 3.1]{IUT1} --- an elliptic curve in Legendre form over
$\mathbb{Q}(\sqrt5)$, with $\ell=7$ --- whose defining prime $p=211$
\emph{splits} in the field of moduli, so that one place of multiplicative
reduction and one place of good reduction lie above it.  All conditions (a)--(f) of \cite[Definition 3.1]{IUT1}
are verified by explicit arithmetic and cited geometric constructions,
independently of the hypothesis used in the comparison below.  For this datum the local invariants are
$e_{\mathrm{g}}=1$, $e_{\mathrm{b}}=105$ and
$\ord_{211}(\underline{\underline q}_{\mathrm{b}})=29/7$, and the extension is
tame at $211$ since $105\le 209=p-2$.

We then examine what this datum implies for the local estimates of
\cite[Theorem 1.10, proof, Step (v)]{IUT4}.  Under one hypothesis about the source, which we
isolate as \SA{} and state precisely in \S\ref{sec:sa}, the log-volume of the
holomorphic hull of the union of the possible images of a $\Theta$-pilot object
at a mixed tuple is bounded below by $-106/105$ in the relevant normalization,
whereas the upper bound obtained by applying \cite[Proposition 1.4(iii)]{IUT4} to
the fixed target component is $-226/105$.  The difference is $8/7>0$.

We are explicit about the status of every step.  The construction of
\S\ref{sec:datum}--\S\ref{sec:profile} is interpretation-free.  The hypothesis
\SA{} of \S\ref{sec:sa} is a reading of the source, and we say exactly which
sentences it rests on and exactly what would falsify it.  We do \emph{not} claim
that the $abc$ conjecture is false, that the lower bound of
\cite[Corollary 3.12]{IUT3} is false, that the global inequality of
\cite[Theorem 1.10]{IUT4} is false, or that the objections of \cite{SS} are
thereby vindicated.
\end{abstract}

\tableofcontents

%=====================================================================
\section{Introduction}\label{sec:intro}

\subsection{What this paper does}

Inter-universal Teichm\"uller theory \cite{IUT1,IUT2,IUT3,IUT4} begins from a
collection of data called an \emph{initial $\Theta$-datum}
\cite[Definition 3.1]{IUT1}.  Explicit initial $\Theta$-data are needed whenever
one wants numerically effective consequences; \cite{MFHMP} is devoted to
obtaining such consequences.  The present paper contributes an explicit initial
$\Theta$-datum with the following feature: the rational prime underlying the bad place \emph{splits} in
the field of moduli, so that the set $\VV$ of \cite[Definition 3.1(e)]{IUT1} 
contains \emph{two} places above it, one bad and one good.

We call this a \emph{mixed profile}.  The number of selected places is fixed
by \cite[Definition 3.1(e)]{IUT1}, which requires $\VV$
to be a \emph{section} of $\VV(K)\twoheadrightarrow\Vmod$, so that the number of
places of $\VV$ above a rational prime $p$ equals the number of places of $\Fmod$
above $p$.  In the present quadratic field, a mixed profile at $p$ occurs when $p$ splits
and the reduction types differ at its two places.

Part~\ref{part:datum} constructs such a datum and verifies
\cite[Definition 3.1]{IUT1} in full.  Part~\ref{part:estimates} examines the
consequences for the local estimates of \cite[Theorem 1.10, proof, Step (v)]{IUT4}.

\subsection{The two results}\label{sec:results}

\begin{theorem}[the datum; \S\ref{sec:datum}--\S\ref{sec:profile}]\label{thm:A}
Let
\[
 F_0=\mathbb{Q}(\sqrt5),\quad \pi=16+3\sqrt5,\quad a=\pi^{29},\quad
 E_0:y^2=x(x-1)(x-a),\quad \ell=7,
\]
\[
 F=F_0\bigl(\sqrt{-1},E_0[15]\bigr),\qquad K=F\bigl(E_0[7]\bigr).
\]
Let $\Vbad_{\mathrm{mod}}$ consist of all places of $F_0$ of multiplicative
reduction outside residue characteristics $2,7$; it contains $v_\pi$, where
$v_\pi(\pi)=1$.  There exist a global cover $\underline C_K$, a compatible
section $\VV$ of $\VV(K)\twoheadrightarrow\Vmod$, and a cusp $\epsilon$ such that
these data form an initial $\Theta$-datum in the sense of
\cite[Definition 3.1]{IUT1}.  The choices are constructed in
\S\ref{sec:e}; an arbitrary section is not asserted to be compatible.  Above the rational prime $211=\Norm(\pi)$ the set
$\VV$ has exactly two elements, one of multiplicative and one of good reduction,
with
\[
 e_{\mathrm{g}}=1,\qquad e_{\mathrm{b}}=105,\qquad
 \ord_{211}\bigl(\underline{\underline q}_{\mathrm{b}}\bigr)=\tfrac{29}{7},
\]
and the profile is tame at $211$, with the stronger bound $105\le209=p-2$.
\end{theorem}

Theorem~\ref{thm:A} uses exact computations, standard facts about elliptic
curves and the cited constructions, checking \cite[Definition 3.1]{IUT1} one
condition at a time.  It is independent of \SA{}.

\begin{theorem}[the comparison; \S\ref{sec:local}]\label{thm:B}
Assume the hypothesis \SA{} of \S\ref{sec:sa}.  For the datum of
Theorem~\ref{thm:A}, at stage $j=1$ and the mixed tuple
$t=(\mathrm{g},\mathrm{b})$, the upper bound obtained by applying
\cite[Proposition 1.4(iii)]{IUT4} to the fixed target component is
$-226/105$ in the normalization by $\log211$, while the holomorphic hull of the actual comparison region contains
a lattice whose hull has log-volume $-106/105$.  The two are
incompatible, with margin
\[
 -\frac{106}{105}-\Bigl(-\frac{226}{105}\Bigr)=\frac{8}{7}>0 .
\]
\end{theorem}

Theorem~\ref{thm:B} rests on exactly one hypothesis about the source, isolated as
\SA{} in \S\ref{sec:sa}.  \SA{} is a \emph{reading}, not a computation; we give
its four supporting steps separately in \S\ref{sec:sabasis} and state in
\S\ref{sec:falsify} the four ways it could fail.  Everything else in
Theorem~\ref{thm:B} uses the cited local estimates and elementary lattice
and arithmetic calculations.

\subsection{What this paper does not claim}\label{sec:notclaim}

We state this at the outset because the subject invites misreading.
This paper does \emph{not} claim:
\begin{enumerate}[label=(\roman*)]
\item that the $abc$ conjecture is false, or that any Diophantine consequence
      asserted in \cite{IUT4} is false;
\item that the lower bound of \cite[Corollary 3.12]{IUT3} is false;
\item that the global inequality of \cite[Theorem 1.10]{IUT4} is false;
\item that the objections raised in \cite{SS} are correct, or that the responses
      of \cite{Rpt2018} are incorrect;
\item that the estimates examined here cannot be repaired.
\end{enumerate}
What is asserted, and only under the hypothesis \SA{} of \S\ref{sec:sa}, is that
one specific step --- the application of \cite[Proposition 1.4(iii)]{IUT4} to a
fixed target component of the union of possible images --- does not yield the
bound it is used to yield, for the datum constructed here.

\subsection{The status of each step}

We use three labels throughout.
\begin{description}[leftmargin=2.2em,style=nextline]
\item[{[S]}] A source statement, cited by number.  Quotation marks delimit
  verbatim wording; unquoted descriptions are paraphrases.  Ordinary deductions
  from these statements are written as proofs.
\item[{[P]}] A proved deduction or an exact computation.  Every computation in Part~\ref{part:datum}
  is carried out in $\mathbb{Z}[\sqrt5]$ or in $\mathbb{Z}$ with exact integer
  arithmetic.  Decimal evaluations in \S\ref{sec:averaged} are illustrative;
  the inequalities there are also proved using rational bounds.
\item[{[H]}] A \emph{reading} of the source: a step whose justification is an
  interpretation of the text rather than a computation or a quoted statement.
  There is exactly one such step of substance, namely \SA{} of \S\ref{sec:sa},
  and \S\ref{sec:falsify} says what would falsify it.
\end{description}

\subsection{Sources and versions}

Statement numbers in this paper refer to the author versions of
\cite{IUT1,IUT2,IUT3,IUT4} dated May 2020 (I), December 2020 (II), May 2020
(III) and April 2020 (IV).  We cite by statement number throughout.
The publication metadata in \cite{IUTpub} identify the corresponding journal
articles; we do not assert textual identity between those editions and the
author versions consulted.

%=====================================================================
\part{An explicit initial \texorpdfstring{$\Theta$}{Theta}-datum}\label{part:datum}

\section{The datum}\label{sec:datum}

\begin{definition}\label{def:datum}
Set
\[
  F_0=\mathbb{Q}(\sqrt5),\qquad
  \pi=16+3\sqrt5,\qquad
  a=\pi^{29},\qquad
  E_0:\; y^2=x(x-1)(x-a)\ \text{over } F_0,
\]
\[
  \ell=7,\qquad
  F=F_0\bigl(\sqrt{-1},\,E_0[15]\bigr),\qquad
  K=F\bigl(E_0[7]\bigr),
\]
and let $X_F\subset E_F:=E_0\times_{F_0}F$ be the once-punctured elliptic curve
obtained by removing the origin.  Let $v_\pi,v_{\bar\pi}$ denote the two places
of $F_0$ above $211$ (see Lemma~\ref{lem:split}), let
$\Vbad_{\mathrm{mod}}$ be the set of all multiplicative places of $F_0$
outside residue characteristics $2,7$.  Fix a line $W\subset E_0[7](K)$ and a
nonzero class $\xi\in E_0[7](K)/W$.  Equip these data with the global cover,
cusp and compatible section $\VV$ constructed in \S\ref{sec:e}.
Thus the arithmetic input is explicit; its auxiliary geometric choices are
specified up to the compatible choices allowed by the source.
\end{definition}

The choices serve the following conditions of
\cite[Definition 3.1]{IUT1}, as follows.
\begin{center}
\begin{tabular}{p{.34\linewidth}p{.60\linewidth}}
\hline
choice & condition it serves\\
\hline
Legendre form $y^2=x(x-1)(x-a)$ & full $2$-torsion over $F_0$, for (b)\\
$E_0[15]\subset E_0(F)$ & semistability everywhere, and $[F{:}\Fmod]$ prime to $7$, for (b)\\
$a=\pi^{29}$, $29$ prime & $\ell\nmid\ord(q)$, for (c)\\
$\Norm(\pi)=211$ split & two places of $\VV$ above $211$, for the mixed profile\\
\hline
\end{tabular}
\end{center}

\section{Elementary arithmetic of the datum}\label{sec:arith}

Write $x\mapsto\bar x$ for the nontrivial automorphism of $F_0$ and
$\Norm=\Norm_{F_0/\mathbb{Q}}$.

\begin{lemma}[P]\label{lem:split}
$\Norm(\pi)=16^2-5\cdot3^2=211$, which is prime.  The congruence $s^2\equiv5
\pmod{211}$ has the two solutions $s\equiv65,146$, so $211$ splits in $F_0$;
write $v_\pi,v_{\bar\pi}$ for the two places, normalized so that
$v_\pi(\pi)=1$, $v_\pi(\bar\pi)=0$.
\end{lemma}

\begin{lemma}[P]\label{lem:val}
Write $a=A+B\sqrt5$ with $A,B\in\mathbb{Z}$ (so $A$ has $40$ decimal digits and
$B$ has $39$).  Then
\[
  \Norm(a)=211^{29},\qquad v_\pi(a)=29,\qquad v_{\bar\pi}(a)=0,
\]
and $211\nmid\Norm(1-a)$, so $v_\pi(a-1)=v_{\bar\pi}(a-1)=0$.
Reducing modulo the two places, $a\equiv0$, $a-1\equiv210$ at $v_\pi$, and
$a\equiv26$, $a-1\equiv25$ at $v_{\bar\pi}$.
\end{lemma}

\begin{proposition}[P]\label{prop:red}
$E_0$ has multiplicative reduction at $v_\pi$ and good reduction at
$v_{\bar\pi}$.  Moreover
\[
 v_\pi(j)=-58,\qquad v_{\bar\pi}(j)=0 .
\]
\end{proposition}

\begin{proof}
The discriminant of the Legendre model is $\Delta=16\,a^2(a-1)^2$ and the
$j$-invariant is $j=256\,(a^2-a+1)^3/\bigl(a^2(a-1)^2\bigr)$.  At $v_\pi$ we have
$v_\pi(a)=29>0$, hence $v_\pi(a^2-a+1)=0$ and $v_\pi(a-1)=0$, so
$v_\pi(\Delta)=2\cdot29=58>0$ and $v_\pi(j)=-2\cdot29=-58<0$; since $211$ is odd,
$v_\pi(16)=0$.  The invariant $c_4=16(a^2-a+1)$ is a unit, so this
integral model has nodal, hence multiplicative, reduction.  At $v_{\bar\pi}$
both $a$ and $a-1$ are units by Lemma~\ref{lem:val}, so
$v_{\bar\pi}(\Delta)=0$; moreover $a^2-a+1\equiv18\pmod{v_{\bar\pi}}$,
which gives $v_{\bar\pi}(j)=0$.
\end{proof}

\begin{proposition}[P]\label{prop:fmod}
$j(E_0)\notin\mathbb{Q}$; consequently the field of moduli of $X_F$ is
$\Fmod=\mathbb{Q}(j)=F_0=\mathbb{Q}(\sqrt5)$.
\end{proposition}

\begin{proof}
The unequal valuations $v_\pi(j)=-58$ and $v_{\bar\pi}(j)=0$ already exclude
$j\in\mathbb Q$, since a rational number has equal valuations at the two split
places.  As an independent exact check, for the Legendre parameter $\lambda=a$
one has $j\in\mathbb{Q}$ if and only if
$j=\bar j$, i.e.\ if and only if $\bar a$ lies in the $S_3$-orbit
$\{\lambda,\lambda^{-1},1-\lambda,(1-\lambda)^{-1},\lambda(\lambda-1)^{-1},
(\lambda-1)\lambda^{-1}\}$ of $a$.  Each of the six equalities is an identity in
$\mathbb{Z}[\sqrt5]$ after clearing denominators, and each fails; the
computation is exact.  Since $j\in F_0$ and $[F_0{:}\mathbb{Q}]=2$, we get
$\mathbb{Q}(j)=F_0$.
\end{proof}

\begin{remark}
Proposition~\ref{prop:fmod} is what makes the mixed profile possible.  Had $j$
been rational we would have had $\Fmod=\mathbb{Q}$, hence a single place of
$\VV$ above $211$ by \cite[Definition 3.1(e)]{IUT1}, and no mixed tuple.
\end{remark}

\section{Verification of the initial-data conditions}
\label{sec:def31}

\subsection{(a): the base field}
Fix an algebraic closure $\overline F$ containing all the fields in
Definition~\ref{def:datum}. The field $F$ is a number field and contains $i$.
This verifies \cite[Definition 3.1(a)]{IUT1}.

\subsection{(b): reduction, the selected places, and the base extension}

\begin{proposition}\label{prop:b}
The curve $X_F$ and the fields of Definition~\ref{def:datum} satisfy
\cite[Definition 3.1(b)]{IUT1}, with
\[
 \Vbad_{\mathrm{mod}}
 =\{u\in\Vmod^{\mathrm{non}}:u\nmid14,\ \ord_u(a(1-a))>0\}.
\]
\end{proposition}

\begin{proof}
For a finite place $w$ of $F$, choose $r=3$ if the residue characteristic is
not $3$, and $r=5$ otherwise. All $r$-torsion points of $E_F$ are rational
over $F_w$, and $r\ge3$ is invertible in $\mathcal O_{F_w}$.
By \cite[Proposition 1.8(v)]{IUT4}, $E_F$ has semistable reduction over
$\mathcal O_{F_w}$ itself. With the origin as marked section this gives
stable reduction of the once-punctured curve $X_F$.

At an odd prime ideal dividing $a(1-a)$, precisely one of $a,1-a$ has
positive order. The invariant $c_4=16(a^2-a+1)$ is a unit there, so the
reduction is multiplicative. When $a\equiv0$ the nodal tangent slopes
are $\pm i$, and when $a\equiv1$ they are $\pm1$.
Consequently the reduction is split multiplicative over $F$, which contains
$i$. The selected set is finite and nonempty, since it contains $v_\pi$;
its residue characteristics are odd and different from $7$.

By Proposition~\ref{prop:fmod}, $\Fmod=F_0$. The field
$F=F_0(i,E_0[15])$ is Galois over $F_0$, and
\[
 [F:F_0]\mid 2\,|\mathrm{GL}_2(\mathbb Z/15\mathbb Z)|
 =2\cdot48\cdot480=46080=2^{10}3^2 5.
\]
This degree is prime to $7$.
Finally $E_0[2]$ is rational over $F_0$ by the Legendre equation, and
$E_0[3]$ is rational over $F$. Thus all $6$-torsion points are rational over $F$.
\end{proof}

\subsection{(c): the torsion image and all selected local orders}

\begin{lemma}[descent]\label{lem:descent}
If $\rho_7(G_{F_0})$ contains $\mathrm{SL}_2(\mathbb F_7)$, then so does
$\rho_7(G_F)$.
\end{lemma}

\begin{proof}
The latter image is normal in the former, and the quotient has order dividing
$[F:F_0]$, hence prime to $7$. Every element of order $7$ of
$\mathrm{SL}_2(\mathbb F_7)$ therefore belongs to $\rho_7(G_F)$.
The upper and lower unipotent root subgroups generate
$\mathrm{SL}_2(\mathbb F_7)$, proving the assertion.
\end{proof}

\begin{proposition}\label{prop:sl2}
The image $\rho_7(G_F)$ contains $\mathrm{SL}_2(\mathbb F_7)$.
\end{proposition}

\begin{proof}
At $v_\pi$, the minimal discriminant order is $58$.
After the unramified splitting extension $\mathbb Q_{211}(i)$,
Tate uniformization identifies the torsion with
$\mu_7$ and a class of $q^{1/7}$.
Because $7\nmid58$ and $211\ne7$, inertia contains a nonidentity
transvection $T$ of order $7$. This conclusion is unchanged by the
unramified splitting extension.

At the place above $11$ given by $\sqrt5\equiv4$, the parameter $a$ reduces
to $2$, and the exact point count is $\#E_0(\mathbb F_{11})=12$.
Thus a Frobenius element $h$ has characteristic polynomial $X^2+11$.
Its discriminant modulo $7$ is $5$, a nonsquare, so $h$ fixes no line
in $\mathbb F_7^2$. The transvections $T$ and $hTh^{-1}$ have different
fixed lines. In a basis given by these two lines they lie in the two
opposite root subgroups. Their powers give both full root subgroups,
since the coefficient field is the prime field $\mathbb F_7$.
Hence $\rho_7(G_{F_0})$ contains $\mathrm{SL}_2(\mathbb F_7)$.
Apply Lemma~\ref{lem:descent}.
\end{proof}

\begin{proposition}\label{prop:coprime}
At every place of $\VV(F)^{\mathrm{bad}}$, the order of the Tate parameter
of $E_F$ is prime to $7$.
\end{proposition}

\begin{proof}
The exact norm certificate supplied with the computation gives
\[
 \Norm(a)=211^{29},\qquad
 \Norm(1-a)=2^2 3^2\cdot5\cdot59\cdot80621\cdot C,
\]
where
\[
 C=29622281599338016545834674091843816951690863770835007573979,
\]
and the second norm is seventh-power-free. The certificate checks all
primes up to $\lfloor C^{1/7}\rfloor=225469788$; it does not assert
that $C$ is prime.
For integral $z\ne0$, the identity
\[
 \ord_q\Norm(z)=\sum_{\mathfrak p\mid q}
 f(\mathfrak p/q)\ord_{\mathfrak p}(z)
\]
shows that every positive order of $1-a$ is at most $6$.
The only positive order of $a$ is $29$ at $v_\pi$.

For a selected prime ideal $\mathfrak p$ of $F_0$, let $r$ be the positive
order of $a$ or $1-a$. Since $c_4$ is a unit, the minimal discriminant order
and Tate parameter order over $F_{0,\mathfrak p}$ are $2r$.
Thus $7\nmid2r$.
At a place $w$ of $F$ above $\mathfrak p$ this order becomes
$2r\,e(w/\mathfrak p)$. Since the ramification index divides the
prime-to-$7$ degree $[F:F_0]$, the resulting order is also prime to $7$.
\end{proof}

Together with Proposition~\ref{prop:sl2}, this verifies all parts of
\cite[Definition 3.1(c)]{IUT1}: $\ell=7$ is prime, $\ell\ge5$, and no
selected residue characteristic is $7$.
Moreover $K=F(E_0[7])$ is exactly the field fixed by the kernel of this
representation, rather than an arbitrary larger torsion-containing field.

\subsection{(d): a fixed global quotient and its specified core}
\label{sec:kcore}

Fix a line $W\subset V:=E_0[7](K)$ and a nonzero element $\xi\in Q:=V/W$.
Let $\pi_W:E_K\to E':=E_K/W$, and let $\psi:E'\to E_K$ be the isogeny
satisfying $\psi\pi_W=[7]$. The group
$\ker\psi=\pi_W(E_K[7])\simeq Q$ is constant over $K$.
Put
\[
 \underline X_K=E'\setminus\ker\psi,\qquad
 \underline C_K=[\underline X_K/\{\pm1\}],\qquad
 C_K=[X_K/\{\pm1\}].
\]
The degree-$7$ covering $\underline C_K\to C_K$ has type
$(1,7\text{-tors})^{\pm}$ in the sense of
\cite[Definition 2.1]{EtTh}.
Indeed, the isogeny has constant kernel, and its zero cusp is rational.
The resulting arithmetic quotient onto $Q$ is surjective on the geometric
group and zero on the decomposition group of the zero cusp, as required
by that definition. The quotient-stack square with
$\underline X_K\to X_K$ is cartesian and finite etale, and hence gives the
corresponding open inclusions of fundamental groups in
\cite[Definition 3.1(d)]{IUT1}.
Let $\epsilon$ be the cusp represented by the sign class of the point
corresponding to $\xi$ in $\ker\psi$.

\begin{proposition}\label{prop:kcore}
The $K$-core of $\underline C_K$ is the specified orbicurve $C_K$.
The corresponding statement remains true over each completion $K_v$.
\end{proposition}

\begin{proof}
By \cite[Proposition 2.1]{MFHMP} (see also \cite[Table 4]{Sijs}), an arithmetic once-punctured elliptic
curve has one of the following four $j$-invariants:
\[
 \frac{488095744}{125}=2^{14}31^3 5^{-3},\qquad
 \frac{1556068}{81}=2^2 73^3 3^{-4},\qquad 1728,\qquad0.
\]
All four are rational, whereas $j(E_0)$ is not by
Proposition~\ref{prop:fmod}.
Thus $X_{\overline K}$ is non-arithmetic; the finite etale commensurable
orbicurve $C_{\overline K}$ is non-arithmetic as well.
It is a punctured hemi-elliptic orbicurve, so
\cite[Proposition 2.7]{CanLift} says that $C_{\overline K}$ is itself a core.
By \cite[Proposition 2.3(i)]{CanLift}, this core descends to $C_K$.
Since $\underline C_K\to C_K$ is finite etale, it has this same core.
For a completion, choose an embedding
$\overline K\hookrightarrow\overline{K_v}$, apply
\cite[Proposition 2.3(ii)]{CanLift}, and then descend by
\cite[Proposition 2.3(i)]{CanLift}. This gives precisely $C_{K_v}$.
\end{proof}

\begin{remark}
Two exceptional $j$-invariants are not algebraic integers.
Non-integrality alone would not exclude them; the preceding argument
uses the nonrationality of $j(E_0)$.
\end{remark}

The geometric covers of types $(1,7\text{-tors}_{\Theta})$ and
$(1,7\text{-tors}_{\Theta})^{\pm}$ recalled in
\cite[Definition 3.1(d)]{IUT1} are the covers supplied by
\cite[Proposition 2.2, Definition 2.3]{EtTh}.
Here $7$ is odd. These geometric covers are distinguished from their
arithmetic models at individual bad places below.

\subsection{(e) and (f): a compatible section and the local covers}
\label{sec:e}

We retain the single global pair $(W,\xi)$, the quotient $\underline C_K$,
and the cusp $\epsilon$ just constructed.
For a selected moduli place $u$, first choose a place $w$ of $F$ over $u$
and a place $v_0$ of $K$ over $w$.
The curve over $F_w$ is split Tate. Its graph quotient on $7$-torsion is
\[
 0\longrightarrow\mu_7\longrightarrow E[7]
 \longrightarrow\mathbb Z/7\mathbb Z\longrightarrow0,
\]
where the class of $q^{1/7}$ maps to $1$.
This determines a local marked flag $(W_{v_0},\pm\xi_{v_0})$.
Changing the root by $\mu_7$ does not change the quotient class.

The group $\mathrm{SL}_2(V)$ acts transitively on pairs consisting of a
line and a nonzero quotient element: lift the quotient element to $t$ and
choose $w$ in the line so that $\det(w,t)=1$, and map one such basis to
another. Since $\mathrm{Gal}(K/F)$ contains this group, choose an element
$\sigma_u$ carrying the local marked flag to $(W,\pm\xi)$, and replace
$v_0$ by $v(u)=\sigma_u v_0$.
With the convention $\iota_{\sigma v}=\iota_v\circ\sigma^{-1}$ for local
embeddings, the local marked flag at $\sigma v$ is the image under $\sigma$
of the one at $v$.
The decomposition group preserves the split Tate graph quotient and its
canonical generator, so this prescription is independent of the representative
of the place.
We are acting on places, not replacing the fixed global quotient or cusp.

In the split Tate model, the quotient map is
$E(q)\to E(q^7)$, induced by $z\mapsto z^7$.
Its dual $\psi:E(q^7)\to E(q)$ is induced by $z\mapsto z$ and has
kernel given by the classes $q^r$ for $r\in\mathbb Z/7\mathbb Z$.
Thus at the chosen place the local covering has type
$(1,\mathbb Z/7\mathbb Z)^{\pm}$, and the cusp $\epsilon_v$ is the canonical
$\pm1$ cusp of \cite[Definition 2.5(i)]{EtTh}.
This verifies the simultaneous local requirements of
\cite[Definition 3.1(e),(f)]{IUT1}.
Choose one lift at every remaining moduli place to complete the section
$\VV\to\Vmod$.
Different moduli places may use different $\sigma_u$; no equivariance of
the section is required. Each $\sigma_u$ fixes $F$, so this construction
preserves all the local ramification indices and parameter orders already
computed.

The local theta models also satisfy the hypotheses used by the source.
The full $2$-torsion gives the condition $K_v=\ddot K_v$ noted in
\cite[Definition 3.1(e)]{IUT1}.
The Weil pairings give $\mu_3\subset F$ and $\mu_7\subset K$, and hence
$\mu_{12}\subset F$ because $i\in F$.
The square and seventh roots of the local Tate parameter exist from the
full $2$- and $7$-torsion.
The local core is non-arithmetic by Proposition~\ref{prop:kcore}.
Use the normalized theta construction of
\cite[Theorem 1.10(iii), Definition 2.5(i)]{EtTh} to give the geometric
theta covers their natural arithmetic models at each bad place.
The auxiliary double cover in that construction can be chosen from a
nonzero rational $2$-torsion element different from the toric one.
No additional field extension or common global arithmetic theta twist
is imposed here.

Finally, construct the arrow-marked covers from $(X_K,\underline C_K,\epsilon)$
as in \cite[Section 1, Definition 1.1, Remark 1.1.2]{IUT1}.
The torsion hypothesis there concerns the original, un-underlined $X_K$;
it follows from $E_0[7](K)$ and does not require all of $E'[7]$ to be rational.
The section used in that construction is supplied by the decomposition group
of the cusp $2\epsilon$; it is available since $7\ge5$ and
$2\xi\ne\pm\xi$.
Base-changing these covers to the good places supplies the remaining local
covers and open subgroups in \cite[Definition 3.1(f)]{IUT1}.
Together with the preceding constructions, this verifies every clause of
\cite[Definition 3.1]{IUT1} for the chosen data.

\begin{proposition}\label{prop:mixed}
The section $\VV$ has exactly two places above $211$, one in $\Vbad$ and
one in $\Vgood$.
\end{proposition}

\begin{proof}
The section is in bijection with $\Vmod=\VV(F_0)$, which has exactly two
places above $211$. Of these, $v_\pi$ belongs to the selected bad set and
$v_{\bar\pi}$ has good reduction. The choices in the preceding construction
preserve these reduction types.
\end{proof}

\subsection{The remaining input for the local estimate}
\label{sec:lattice}

\begin{lemma}\label{lem:theater-existence}
The datum above admits an LGP-Gaussian log-theta-lattice as in
\cite[Definition 3.8(iii)]{IUT3}.
\end{lemma}

\begin{proof}
Start with the native models of the initial data in
\cite[Examples 3.2--3.5, 4.3, 4.4, 5.4, 6.2, 6.3]{IUT1}.
The bridge identifications of
\cite[Proposition 4.8(ii), Proposition 6.7]{IUT1} align the additive and
multiplicative models to give a full theater of
\cite[Definition 6.13]{IUT1}.
The Frobenioid models are constructed before applying the isomorphism
lifting of \cite[Corollary 5.3(ii)]{IUT1}.
On tagged copies indexed by $\mathbb Z^2$, form the synchronized full
log-links of \cite[Proposition 1.3]{IUT3}: one theater isomorphism is used
for all constituents before taking the full collection.
For each horizontal edge, use the nonempty enhanced Gaussian isomorphism
sets of \cite[Corollary 4.10(iii)]{IUT2} and the Gaussian/LGP identifications
relative to the incoming log-link in \cite[Proposition 3.7]{IUT3}.
Taking the full isomorphism sets supplies the links required by
\cite[Definition 3.8(ii),(iii)]{IUT3}.
This construction uses neither a numerical inequality from
\cite[Corollary 3.12]{IUT3} nor the hypothesis \SA{} below.
\end{proof}

For the additional arithmetic hypotheses of \cite[Theorem 1.10]{IUT4},
the Legendre model has $F_{\mathrm{tpd}}=F_0$ and
\[
 F=F_0(i,E_0[30])=F_0(i,E_0[15])
\]
exactly. At every nonselected place outside $2,7$, the discriminant
$16a^2(1-a)^2$ is a unit, so the curve has good reduction.
Thus the extra good-reduction and exact-field hypotheses hold.

\section{The local profile at \texorpdfstring{$211$}{211}}\label{sec:profile}

\begin{proposition}\label{prop:profile}
With notation as above, write $\mathrm{b}$ and $\mathrm{g}$ for the bad and good
places of $\VV$ above $211$.  Then
\[
  e_{\mathrm{g}}=1,\qquad e_{\mathrm{b}}=105,\qquad
  \ord_{211}\bigl(\underline{\underline q}_{\mathrm{b}}\bigr)=\frac{29}{7},
\]
and the extension is tame at $211$, since $211\nmid105$; moreover
$105\le209=p-2$.
\end{proposition}

\begin{proof}
Since $15$ and $7$ are coprime, the definition of $K$ gives the exact equality
\[
 K=F_0(i,E_0[105]).
\]
At either place of $F_0$ above $211$, the completion is $\mathbb Q_{211}$,
and adjoining $i$ is unramified. At the good place,
\cite[Proposition 1.8(vii)]{IUT4} shows that the $105$-torsion field
is unramified, hence $e_{\mathrm g}=1$.

At the bad place, pass first to the unramified field
$k_0=\mathbb Q_{211}(i)$. The curve then has split multiplicative reduction.
We use the standard Tate uniformization fact that the full $n$-torsion field
of a split Tate curve is $k_0(\mu_n,q^{1/n})$.
Here $\ord_{211}(q)=58$ and $n=105$.
The value $58/105$ of $q^{1/105}$, in lowest terms, forces the ramification
index to be divisible by $105$.
On the other hand, \cite[Proposition 1.8(vii)]{IUT4} says that it divides
$105$. Thus $e_{\mathrm b}=105$.
Finally \cite[Example 3.2(iv)]{IUT1} supplies the $2\ell$-th root
\[
 \underline{\underline q}_{\mathrm b}=q^{1/(2\ell)}
 \quad\text{up to a }\mu_{2\ell}\text{-multiple},
 \qquad
 \ord_{211}(\underline{\underline q}_{\mathrm b})=\frac{58}{14}=\frac{29}{7}.
\]

The normalized valuation $\ord_{211}$ always has $\ord_{211}(211)=1$.
In contrast, the discrete orders of $q$ at the selected $F$- and $K$-places
are $870$ and $6090$, respectively; those of the theta root are measured
with the same normalization when compared in Part~\ref{part:estimates}.
\end{proof}

\begin{remark}[why tameness matters]\label{rem:tame}
[S] Put
\[
 L_I:=\bigotimes_{i\in I,\,\mathbb Z_p}\log_p(R_i^\times).
\]
Thus $L_I$ is the object denoted by $\log_p(R_I^\times)$ in
\cite[Proposition 1.2]{IUT4}; this is a tensor-log notation.
The orders of differents are normalized by $\ord_p(p)=1$.
For tame $k_i/\mathbb Q_p$, \cite[Proposition 1.3(i)]{IUT4},
with $k_0=\mathbb Q_p$, gives $d_i=1-1/e_i$.
If moreover $p>2$ and $e_i\le p-2$, then
\cite[Proposition 1.2]{IUT4} gives $a_i=1/e_i=-b_i$.
For the subset $I^*$ in \cite[Proposition 1.4(iii)]{IUT4},
the two displayed upper bounds are
\[
 \mu^{\log}\!\left(p^{\lfloor\lambda-d_I-a_I\rfloor}L_I\right)
 \le\left(-\lambda+d_I+1+\frac{4|I^*|}{p}\right)\log p
\]
and
\[
 \mu^{\log}\!\left(
 p^{\lfloor\lambda-d_I-a_I\rfloor-b_I}(R_I)^{\sim}\right)
 \le(-\lambda+d_I+1)\log p+
 \sum_{i\in I^*}(3+\log e_i).
\]

In \cite[Theorem 1.10, proof, Step (v)]{IUT4}, the distinguished slot
is $i^\dagger=j$, and the substitution in these formulas is
\[
 \lambda=0\quad(t_j\in\Vgood),\qquad
 \lambda=\ord_p(\underline{\underline q}_{t_j}^{\,j^2})
        =j^2\ord_p(\underline{\underline q}_{t_j})\quad(t_j\in\Vbad).
\]
The formulas below with $\lambda=29/7$ concern $j=1$.

The parameter $\lambda$ in these source formulas belongs to
$e_{i^\dagger}^{-1}\mathbb Z$ for some $i^\dagger\in I$; here
$29/7=435/105$ satisfies that condition.
For the profile of Proposition~\ref{prop:profile} every $e_i$ satisfies
$e_i\le105\le209=p-2$, so one may take $I^*=\varnothing$ and \emph{both} error
terms vanish: the bound is exactly $(-\lambda+d_I+1)\log p$.  For this small-ramification profile the cited formulas give
\[
 d_I=0+\tfrac{104}{105}=\tfrac{104}{105},\qquad
 a_I=1+\tfrac1{105}=\tfrac{106}{105},\qquad
 d_I+a_I=2=|I| ,
\]
attaining the source's inequality $d_I+a_I\ge|I|$ with equality; and the exponent
occurring in the intermediate container is exactly $\lfloor\lambda\rfloor-|I|$.
This is why the comparison of \S\ref{sec:local} can be made without error terms.
\end{remark}

%=====================================================================
\part{The local estimates}\label{part:estimates}

\section{The tensor packet and the tuple index}\label{sec:packet}

[S] \cite[Proposition 3.2]{IUT3} defines, for a capsule $\{{}^\alpha D^\vdash\}_{\alpha\in A}$
of $D^\vdash$-prime-strips, the $1$-tensor packet as
$\log({}^\alpha D^\vdash_{v_{\mathbb{Q}}})=\bigoplus_{\VV\ni v\mid v_{\mathbb{Q}}}
\log({}^\alpha D^\vdash_v)$ and the $n$-tensor packet as
$\log({}^A D^\vdash_{v_{\mathbb{Q}}})=\bigotimes_{\alpha\in A}
\log({}^\alpha D^\vdash_{v_{\mathbb{Q}}})$.

\begin{observation}\label{obs:tuple}
Consequently, with $A=S^\pm_{j+1}$ and $v_{\mathbb{Q}}=211$,
\[
 I^{\mathbb{Q}}\bigl(S^\pm_{j+1};{}^{n,\circ}D^\vdash_{211}\bigr)
 \;=\;\bigotimes_{\alpha\in S^\pm_{j+1}}\Bigl(\bigoplus_{v\in\VV,\ v\mid211}\cdots\Bigr)
 \;=\;\bigoplus_{t\,:\,S^\pm_{j+1}\to\VV_{211}}\ \bigotimes_{\alpha}(\cdots)_{t(\alpha)} .
\]
By Proposition~\ref{prop:mixed}, $|\VV_{211}|=2$ with one bad and one good
member, so the summands are indexed by tuples
$t\in\{\mathrm{g},\mathrm{b}\}^{\,j+1}$, of which there are $2^{j+1}$.
The outer tensor index is the \emph{label} set, on which $\Ind1$ acts; the inner
direct-sum index is the set of \emph{places}; $\Ind2$ acts independently
within each place-indexed summand
(\cite[Theorem 3.11(i)]{IUT3} records that the number of direct summands equals
the number of $v\in\VV$ over $v_{\mathbb{Q}}$).
\end{observation}

For $j\ge1$, nonconstant \emph{place} tuples exist precisely when
$|\VV_p|\ge2$, regardless of reduction types.  A tuple containing both
chosen-good and chosen-bad places is a further condition.  Under the additional
good-reduction assumption of \cite[Theorem 1.10]{IUT4}, these chosen types
coincide with the actual reduction types outside $2\ell$.
No statistical independence of bad primes is assumed here.

Observation~\ref{obs:tuple} is what connects Part~\ref{part:datum} to the
estimates: the tuple index of the estimates and the split-prime structure of the
datum are the same object.

\section{The hypothesis \texorpdfstring{\SA}{(SA)}}\label{sec:sa}

Fix one lattice supplied by \S\ref{sec:lattice}, one column $(n,\circ)$,
a stage $j$, and $p=211$.  The fixed ambient is
\[
 E=I^{\mathbb{Q}}(S^\pm_{j+1};{}^{n,\circ}D^\vdash_{p}).
\]
Write $U^{\mathrm{act}}_{j,t}$ for the
$t$-component of the union of possible images of a $\Theta$-pilot object in the
sense of \cite[Corollary 3.12]{IUT3}, $C_{\widetilde O_t}(-)$ for the terminal
holomorphic hull at the fixed arithmetic structure, and
$L_t=\bigotimes_i\log\mathcal{O}_{K_{t_i}}^\times$.

\begin{definition}[\SA]\label{def:sa}
\[
\begin{cases}
 L_t\subseteq C_{\widetilde O_t}\bigl(U^{\mathrm{act}}_{j,t}\bigr),
   & \text{$t$ contains at least one good factor,}\\[2pt]
 p^{\lceil j^2N/7\rceil}L_t\subseteq C_{\widetilde O_t}\bigl(U^{\mathrm{act}}_{j,t}\bigr),
   & t=(\mathrm{b},\dots,\mathrm{b}),
\end{cases}
\]
where $N=\ord_{211}(\underline{\underline q}_{\mathrm{b}})\cdot 7$.
\end{definition}

\subsection{The source basis for \SA{}}\label{sec:sabasis}

[H] We now give the reading.  It has four steps; none of them is a general
theorem about functors, and each rests on a separately identified passage.

\begin{enumerate}[label=\textbf{(SA\arabic*)},leftmargin=3.2em]
\item \emph{One output exists.}  [S] \cite[Proposition 3.7(ii),(v)]{IUT3}
  constructs the local fractional ideal $J_v$ and states that it ``is a subset
  [indeed a submodule when $v\in\VV^{\mathrm{non}}$] of
  $I^{\mathbb{Q}}({}^{A,\alpha}F_v)$ whose $\mathbb{Q}$-span is equal to
  $I^{\mathbb{Q}}({}^{A,\alpha}F_v)$'', and builds the associated global objects;
  \cite[Definition 3.8(i)]{IUT3} defines the $\Theta$-pilot object from these;
  \cite[Proposition 3.10(i)]{IUT3} gives the compatibility of the native and
  vertically coric algorithms with respect to the specified labelled Kummer
  maps; \cite[Theorem 3.11(ii)(a)]{IUT3} gives ``isomorphisms with local
  mono-analytic tensor packets and their $\mathbb{Q}$-spans \ldots\ all of which
  are compatible with the respective log-volumes''.  Here the comparison in
  Theorem 3.11(ii) is taken before imposing $\Ind1,\Ind2$; its additive map
  $\kappa_E$ is distinct from the map on Frobenioid objects.  The labelled
  compatibility in Proposition 3.10(i) connects these two realizations.
  Under our reading, realizing the specified native representation yields one region
  $R=\kappa_E[J]$ which is among the possible images of
  \cite[Corollary 3.12]{IUT3}.  We do \emph{not} claim that every isomorphic
  $fJ$, or the full $\mathbb{Q}$-span, is a possible image.
\item \emph{Which automorphisms are admitted.}  [S] In \cite[\S0 and
  Definition 4.10]{IUT1} the arrows of a procession are not single fixed
  inclusions but sets of capsule-full poly-morphisms.  Hence the factor
  transposition required below is realized by a single automorphism $\sigma$ of
  the whole procession.  Preserving the one-element capsule
  ${}^{n,\circ}D^\vdash_0$ is not a condition fixing label $0$
  inside capsules with two or more factors; the slot distinguished by Step (v)
  of \cite[Theorem 1.10]{IUT4} is $j$.
\item \emph{The type of the action.}  [S] \cite[Theorem 3.11(i)]{IUT3} treats the
  collection of data (a),(b),(c) ``regarded up to indeterminacies'' of which
  $\Ind1$ is ``the indeterminacies induced by the automorphisms of the procession
  of $D^\vdash$-prime-strips $\mathrm{Prc}({}^{n,\circ}D^\vdash_T)$'', and states
  that the algorithm is ``functorial with respect to isomorphisms of processions
  of $D^\vdash$-prime-strips''.  An automorphism is an isomorphism, so one
  obtains an action $T_\sigma:E\to E$ transporting labels, coefficients and the
  local and global realizations \emph{together}.  In particular one does not move
  coefficients alone, nor recombine components of distinct tuples.
\item \emph{Membership in the same comparison target.}  [S]
  \cite[Corollary 3.12]{IUT3} defines $U_{j,p}\subseteq E$ as the \emph{union} of
  the possible images ``relative to the relevant Kummer isomorphisms \ldots\
  which we regard as subject to the indeterminacies $\Ind1$, $\Ind2$, $\Ind3$'',
  and only then forms the holomorphic hull.  The introduction of \cite{IUT3}
  states the same thing: the quantity is the log-volume of the holomorphic hull
  ``of the union of the collection of possible images \ldots\ where we recall
  that these `possible images' are subject to the indeterminacies $\Ind1$,
  $\Ind2$, $\Ind3$''.  Applying (SA1)--(SA3) to this definition places both $R$
  and $T_\sigma R$ in the same $U_{j,p}$.  The restriction in the proof of
  \cite[Corollary 3.12]{IUT3} to possible images interpretable as an
  $\mathcal{F}^\Vdash$-prime-strip is preserved, since the transport carries the
  localizations and the global object together.
\end{enumerate}

\begin{proof}[Derivation of \SA{} from (SA1)--(SA4)]
Suppose $t$ contains a good factor.  Choose $\sigma$ so that
$s=\sigma^{-1}t$ has a good factor in its distinguished slot.  The insert-$1$
construction of \cite[Propositions 3.1(ii), 3.4(i),(ii)]{IUT3} and the generated
local fractional ideals of \cite[Example 3.6(ii), Proposition 3.7(v),
Definition 3.8(i)]{IUT3} give $J_s=\widetilde O_s$ for this specified native
representation.  Here $\widetilde O_s$ is the product of the integer rings in
the native tensor algebra's field components, not necessarily the tensor
order $R_s=\bigotimes_i\mathcal O_{K_{s_i}}$.  Thus
$L_s\subseteq R_s\subseteq J_s$.  This is not a claim about every possible image.

The comparison of \cite[Theorem 3.11(ii)(a), Proposition 3.2(i),(ii)]{IUT3}
is an isomorphism of ind-topological modules, including their integral
structures $I$ and their $\mathbb{Q}$-spans.  On the finite-dimensional
portions used here it is additive and continuous, hence $\mathbb{Q}_p$-linear,
and it preserves $I$ as a lattice.  The odd-prime identity
$I_i=p^{-1}L_i$ of \cite[Remark 1.2.2(i)]{IUT3} gives
$L_s=p^{j+1}I_s$ for the integral tensor construction.  Hence this comparison
also carries $L_s$ to the corresponding lattice.  The equality
$L_i=\mathfrak m_i$, used here with $e_i\le105\le p-2$, comes separately from
\cite[Proposition 1.2(i)]{IUT4}.

Transporting by
$T_\sigma$ and taking the $\mathbb{Z}_p$-span and then the terminal hull at the
target gives $L_t\subseteq C_{\widetilde O_t}(U^{\mathrm{act}}_{j,t})$.  For
$t=(\mathrm{b},\dots,\mathrm{b})$, put $c=\lceil j^2N/7\rceil$; the specified
coefficient $a_s$ has valuation $j^2N/7$, so $p^c/a_s$ is integral and
$p^cR_s\subseteq a_sR_s\subseteq J_s$, whence $p^cL_s\subseteq J_s$; since
$\kappa_E$ commutes with the integer scalar $p^c$ and preserves $L$, the same
conclusion follows.  No commutation of $\kappa_E$ with the hull, and no ring
linearity of $\kappa_E$, is used.
\end{proof}

The tuple transport used here satisfies
\[
 \operatorname{pr}_tT_\sigma
 =\tau_{\sigma,t}\operatorname{pr}_{\sigma^{-1}t},\qquad
 \tau_{\sigma,t}:E_{\sigma^{-1}t}\longrightarrow E_t.
\]
It is a tensor-factor permutation and does not require a field isomorphism
$K_{\mathrm b}\simeq K_{\mathrm g}$.  All four steps concern one whole
representation in the same fixed ambient.  A general theorem that functoriality
alone saturates a collection of outputs is neither used nor asserted.

\subsection{Two readings that must be distinguished}\label{sec:twoparts}

The following two statements about $\Ind1$ are \emph{not} equivalent, and the
argument above uses the second, not the first:
\begin{description}[leftmargin=2.4em,style=nextline]
\item[(A) Invariance] the log-volumes are invariant under $\Ind1$
  (\cite[Proposition 3.9(i)--(iii)]{IUT3}; step (x) of the proof of
  \cite[Corollary 3.12]{IUT3});
\item[(B) Closure] the collection of possible images is stable under $\Ind1$
  (\cite[Corollary 3.12]{IUT3} and the introduction of \cite{IUT3}, quoted in
  (SA4)).
\end{description}
(A) does not imply (B); and even (A) together with (B) does not imply that the
hull of the union is bounded by the hull of a member, because the hull of a
union can strictly exceed the hull of every member.  We record this in
Appendix~\ref{app:lean} as a machine-checked finite model, and note that the
phenomenon is present in the source itself in the worked example of
\cite[Remark 3.9.5(iv)]{IUT3}.

\section{The local comparison}\label{sec:local}

\begin{theorem}\label{thm:local}
Assume \SA{}.  For the datum of Definition~\ref{def:datum}, take $j=1$ and the
mixed tuple $t=(\mathrm{g},\mathrm{b})$.  Then, in the normalization by
$\log 211$:
\[
 \text{upper bound from \cite[Proposition 1.4(iii)]{IUT4} with $I^*=\varnothing$: }\;
 -\lambda+d_I+1=-\frac{29}{7}+\frac{104}{105}+1=-\frac{226}{105},
\]
\[
 \text{value of the hull of the seed }L_t:\; -\frac{106}{105},
 \qquad\text{difference }\;\frac{8}{7}>0 .
\]
Here $j=1$, so $\lambda=j^2\ord_{211}(\underline{\underline q}_{\mathrm{b}})=29/7$, and
$d_I=104/105$ by Remark~\ref{rem:tame}, and $I^*=\varnothing$ is admissible
because every $e_i\le105\le p-2$, so no error term occurs.
Consequently the bound is not satisfied by the region it is applied to.
Moreover the stronger container $p^2L_t$ does not contain $L_t$, since
$L\not\subseteq p^2L$ for a nonzero finitely generated
$\mathbb{Z}_{211}$-lattice.
\end{theorem}

\begin{proof}
Put $R_t=\bigotimes_i\mathcal O_{K_{t_i}}$ and choose uniformizers
$\pi_i$ in the factors.  Here the fixed holomorphic integer ring is
$\widetilde O_t=\widetilde R_t$.  Proposition~\ref{prop:profile} and
\cite[Proposition 1.2(i)]{IUT4} give
$L_t=(\bigotimes_i\pi_i)R_t$.  Its ideal hull is
$(\bigotimes_i\pi_i)\widetilde R_t$, whose degree-normalized log-volume
is $-(1+1/105)\log211$ by \cite[Proposition 1.4(i)]{IUT4}.
The containment in \SA{} therefore gives this lower bound for the actual
terminal hull.  The proposed application at the fixed target has the
right-hand side $(-29/7+104/105+1)\log211$, so the asserted comparison
fails by the displayed positive difference.
The issue persists at the level of hulls: $L_t$ is not contained in
$p^2C_{\widetilde O_t}(L_t)$, since its generating tensor has order
$1+1/105$, while multiplication by $p^2$ adds $2$ in every field component.
Thus taking a terminal hull does not rescue that particular container.
In this comparison $C_{\widetilde O_t}(p^2L_t)=p^2C_{\widetilde O_t}(L_t)$;
no commutation of a Kummer map with the hull is needed.
\end{proof}

\begin{remark}
Theorem~\ref{thm:local} does not refute \cite[Proposition 1.4]{IUT4} itself; the
proposition's estimate for the input regions it is stated for is not in question.
What is in question is the application of that estimate to the fixed target
component of the union of possible images.
\end{remark}

\begin{remark}[the nonpositive integer power]
\cite[Theorem 1.10, proof, Step (v)]{IUT4} permits a suitable nonpositive integer power in the
intermediate container.  Enlarging the container is possible, but the resulting
loss of volume must then be carried into the subsequent bound; adding an
$N$-independent constant does not absorb a difference in the slope in $N$.  For
the tame profile of Proposition~\ref{prop:profile} the exponent is exactly
$\lfloor\lambda\rfloor-|I|$, and for $j=1$ mixed this is
$\lfloor29/7-2\rfloor=2$.
\end{remark}

\section{The averaged comparison}\label{sec:averaged}

The explicit datum of Definition~\ref{def:datum} has $N=29$.
In this section we use a different family: for each fixed positive integer
$N\equiv1\pmod{105}$ we construct infinitely many integers $A$ with
$a=A-\sqrt5$.  In particular, $N=211$ is an existence result for this second
family; it does not change the exponent of the explicit datum.
The arithmetic existence argument below does not assume \SA{}.

\begin{proposition}[weights and the conditional lower estimate]\label{prop:weights}
Suppose the two selected places above $p=211$ have
$e_{\mathrm g}=1$, $e_{\mathrm b}=105$ and
$[(\Fmod)_{\mathrm g}:\mathbb Q_p]=[(\Fmod)_{\mathrm b}:\mathbb Q_p]=1$.
At stage $j$, put $k=j+1$ and let $t\in\{\mathrm g,\mathrm b\}^k$.
For the degree-normalized local log-volumes, the weight of $t$ is $2^{-k}$.
Assuming \SA{}, for $N\equiv1\pmod7$ one has
\[
 \frac{h_{211}}{\log211}\ge
 -\frac{53}{35}-\frac{N}{16}-\frac{5}{48}.
\]
\end{proposition}
\begin{proof}
The raw weights in \cite[Remark 3.1.1(ii)]{IUT3} include the factor
$\prod_i[K_{t_i}:(\Fmod)_{t_i}]^{-1}$.  Passing to the degree-normalized
log-volumes of \cite[Proposition 1.4(i)]{IUT4} cancels the local degrees, as
explained explicitly in \cite[Remark 1.7.1]{IUT4}.  The remaining weight is
\[
 \frac{\prod_i[(\Fmod)_{t_i}:\mathbb Q_p]}
 {\sum_{u\in\{\mathrm g,\mathrm b\}^k}\prod_i[(\Fmod)_{u_i}:\mathbb Q_p]}
 =2^{-k}.
\]
No residue degree of $K_{t_i}$ is assumed to equal one here.
By the logarithm calculation with $e_i\le105\le p-2$, the normalized log-volume of the integral
ideal hull of $L_t$ is $-\sum_i e_{t_i}^{-1}\log p$.
Indeed, $L_t=(\bigotimes_i\pi_{t_i})R_t$, and its hull is
$(\bigotimes_i\pi_{t_i})\widetilde R_t$, even when $R_t$ is not normal.
Thus the average of its depth is $53k/105$.
By \SA{}, only the all-bad tuple needs the additional depth
$\lceil j^2N/7\rceil$.  Averaging first over tuples and then over $j=1,2,3$ gives
\[
 \frac{h_{211}}{\log p}\ge
 -\frac13\sum_{j=1}^3\left\{
 \frac{53(j+1)}{105}+2^{-(j+1)}\left\lceil\frac{j^2N}{7}\right\rceil
 \right\}.
\]
For $N\equiv1\pmod7$, the three ceilings are
$(N+6)/7$, $(4N+3)/7$ and $(9N+5)/7$.
The constant, slope and rounding term are respectively
\[
 \frac13\sum_{j=1}^3\frac{53(j+1)}{105}=\frac{53}{35},\quad
 \frac1{21}\left(\frac14+\frac48+\frac9{16}\right)=\frac1{16},\quad
 \frac13\left(\frac6{28}+\frac3{56}+\frac5{112}\right)=\frac5{48}.
\]
\end{proof}

\begin{remark}[the Tate-degree normalization]
For a split semistable elliptic curve over $k$ and $v\mid p$, let $v_v$ be the
discrete valuation, $e_v=v_v(p)$ and $f_v$ the residue degree.
The definition of $D=\deg(q_E)$ in \cite[\S1.2]{SS} gives
\[
 D_p=\frac1{[k:\mathbb Q]}\sum_{v\mid p}v_v(q_v)f_v\log p
 =\frac{\log p}{[k:\mathbb Q]}
   \sum_{v\mid p}[k_v:\mathbb Q_p]\ord_p(q_v).
\]
The summand is zero at good places.
With the full place set pulled back this is invariant under finite extension,
since $\sum_{w\mid v}[k'_w:k_v]=[k':k]$; it is the normalized degree of
\cite[Definition 1.9]{IUT4}.  Here $D_{211}=N\log211$, while the local bad
root has order $N/7$ and its degree contribution is $D_{211}/14$.
The root factor $1/(2\ell)$ and the moduli-place weight $1/2$ are distinct.
The native Tate coefficient in \cite[equation (1.5)]{SS} and the local upper
expression below is $1/3$.  The conditional seed coefficient above is $1/16$;
their difference $13/48$ concerns these two bounds, not the actual volume.
The coefficient in the simplified expression \cite[equation (1.6)]{SS} is
instead $1/14$.  Equality of input normalizations does not identify the
possible-image regions of those arguments.  If the selected bad support is
smaller, this comparison concerns the corresponding restriction of $q_E$.
\end{remark}

\begin{proposition}[the expression in the local upper estimate]\label{prop:B2}
For the family constructed below, denote by $B_2(N)$ the final right-hand
side of \cite[Theorem 1.10, proof, Step (v)]{IUT4}, at $p=211$.
Then
\[
 \frac{B_2(N)}{\log211}=
 \frac{1984}{735}-\frac N3+
 \frac{40}{3}\frac{\log7741440}{\log211}.
\]
This proposition evaluates that expression; it does not assume that it is an
upper bound for the actual family defined using \SA{}.
\end{proposition}
\begin{proof}
The cited final expression is
\[
 \frac{\ell+1}{4}\left\{
 \left(1+\frac4\ell\right)\log\mathfrak d^K_p
 -\frac16\log q_p+\frac4\ell\log\mathfrak s^{\mathbb Q}_p
 +\frac{20}{3}\ell^*_{\mathrm{mod}}\log\mathfrak s^{\le}_p\right\}.
\]
Here $\ell=7$, $e_{\mathrm{mod}}=2$,
$e^*_{\mathrm{mod}}=2^{12}3^3\cdot5\cdot2=1105920$ and
$\ell^*_{\mathrm{mod}}=\log(e^*_{\mathrm{mod}}\ell)=\log7741440$.
The definitions of the support divisors in Step (iii) assign the coefficient
$\iota_p/\log p$ to $p$ in $\mathfrak s^{\le}$.
Since $211$ is ramified in $K$ and is below the cutoff $7741440$,
\[
 \log\mathfrak s^{\le}_{211}=1,\qquad
 \log\mathfrak s^{\mathbb Q}_{211}=\log211.
\]
The two moduli-place weights are $1/2$.  Their normalized different
exponents are $0$ and $104/105$, and their Tate orders are $0$ and $2N$.
Consequently
\[
 \log\mathfrak d^K_{211}=\frac{52}{105}\log211,\qquad
 \log q_{211}=N\log211.
\]
Notice that this last $q$ is the Tate-divisor quantity; the theta parameter
at the bad place has order $N/7$ instead.
Substitution gives the formula, since
$2[(11/7)(52/105)+4/7]=1984/735$.
\end{proof}

\begin{proposition}[an arithmetic existence family]\label{prop:exist}
For every positive integer $N\equiv1\pmod{105}$ there are infinitely many
positive integers $A$ with the following properties.  Put
$a=A-\sqrt5$ and $F_0=\mathbb Q(\sqrt5)$.
At the two places of $F_0$ above $211$, labeled by
$\sqrt5\equiv65$ and $\sqrt5\equiv146$, respectively,
\[
 \ord_{\mathrm b}(a)=N,\qquad
 a\equiv130\pmod{\mathrm g},\qquad
 \ord_{\mathrm b}(1-a)=\ord_{\mathrm g}(1-a)=0.
\]
Moreover,
\[
 \frac{\Norm(a)\Norm(1-a)}{211^N}
\]
is a positive seventh-power-free integer, and
$A\equiv0\pmod8$, $A\equiv2\pmod{3,5,19}$, $A\equiv6\pmod{11}$.
\end{proposition}
The complete congruences and counting proof are given in Appendix~\ref{app:family}.


\paragraph{Arithmetic data for the family.}
For each such $A$ set
\[
 E_0:y^2=x(x-1)(x-a),\quad
 F=F_0(\sqrt{-1},E_0[15]),\quad K=F(E_0[7]).
\]
Choose as $\Vbad_{\mathrm{mod}}$ all multiplicative places of $F_0$ of
residue characteristic other than $2,7$, and choose the quotient, cusp and
section compatibly as in \S\ref{sec:e}.
The same rule for the selected bad set is applied to this new curve.  The arithmetic checks used earlier apply as follows.

The two conjugate $j$-valuations at $211$ are $-2N$ and $0$:
$130^2-130+1\equiv102\not\equiv0\pmod{211}$ on the good side.
Thus $\Fmod=\mathbb Q(j)=F_0$, and the rational exceptional $j$-values in
the core criterion of \S\ref{sec:kcore} are excluded.
For a rational prime $q\ne211$, the norm identity in
Proposition~\ref{prop:exist} implies that every positive local order of
$a$ or $1-a$ is at most six.  At $211$ the only such order is $N$.
At odd multiplicative places the Tate order is twice that order, hence
prime to $7$.  The extension $F/F_0$ is Galois of degree dividing
$2|\operatorname{GL}_2(\mathbb F_3)||\operatorname{GL}_2(\mathbb F_5)|
=46080$, also prime to $7$, so the required coprimality persists over $F$.

At the place of $F_0$ above $11$ with $\sqrt5\equiv4$, the parameter is
$a\equiv2$.  The curve $y^2=x(x-1)(x-2)$ has $12$ points over $\mathbb F_{11}$,
so its Frobenius has trace zero and discriminant $-44\equiv5\pmod7$,
a nonsquare.  The Tate inertia at $211$ supplies a nontrivial transvection.
The same elementary group argument as in \S\ref{sec:def31}, or
\cite[Lemma 3.1(iii)]{GenEll}, then gives
$\operatorname{SL}_2(\mathbb F_7)$ in the image over $F_0$.
The image over $F$ is normal with prime-to-$7$ quotient, so it contains
all the order-$7$ transvections and their generated
$\operatorname{SL}_2(\mathbb F_7)$.
The core and compatible marked-section constructions therefore apply.

At an odd prime the integral Legendre model has good or multiplicative
reduction; its multiplicative reduction is split over $F$, since $i\in F$.
The full $3$- and $5$-torsion and \cite[Proposition 1.8(v)]{IUT4} give
semistability at every finite place.  By the chosen bad set, every good-set
place outside $2\ell$ is genuinely of good reduction, as required by
\cite[Theorem 1.10]{IUT4}.
All $2$-torsion is already $F_0$-rational, so $F$ has exactly the field form
required there, and $d_{\mathrm{mod}}=e_{\mathrm{mod}}=2$.
Finally, the Tate torsion calculation of \S\ref{sec:profile}, using
$\gcd(2N,105)=1$, gives
\[
 e(F_{\mathrm b}/\mathbb Q_{211})=15,\quad
 e(K_{\mathrm b}/\mathbb Q_{211})=105,\quad
 e(K_{\mathrm g}/\mathbb Q_{211})=1,\quad
 \ord_{211}(\underline{\underline q}_{\mathrm b})=\frac{2N}{14}=\frac N7.
\]
These are arithmetic statements and do not use \SA{}; the actual-family
containment needed for the lower estimate remains precisely the separate
hypothesis \SA{}.

\begin{theorem}\label{thm:averaged}
For any datum in the preceding existence family with $N=211$, assume
\SA{}.  Then
\[
 \frac{h_{211}}{\log211}\ge-\frac{12437}{840},\qquad
 \frac{B_2(211)}{\log211}<-\frac{6117}{245},\qquad
 \frac{h_{211}-B_2(211)}{\log211}>\frac{59749}{5880}>0.
\]
In particular, $B_2(211)$ cannot bound that actual $211$-local log-volume.
\end{theorem}
\begin{proof}
The first bound is Proposition~\ref{prop:weights}.
The elementary exponential series gives
$e>65/24$ and
$e<5/2+(1/6)/(1-1/4)=49/18<11/4$.
The integer comparisons
\[
 65^{16}>7741440\cdot24^{16},\qquad 11^5<211\cdot4^5
\]
therefore imply $\log7741440<16$ and $\log211>5$.
Substitute $\log7741440/\log211<16/5$ into
Proposition~\ref{prop:B2}.  The resulting upper bound is $-6117/245$;
subtracting it from $-12437/840$ gives $59749/5880$.
\end{proof}

\begin{remark}
Weakening $5/48$ to $7/48$ gives instead
$h_{211}/\log211\ge-1559/105$ and the still positive margin
$7438/735$.  The two certified margins differ by exactly $1/24$.
The proof uses no floating-point comparison.  This conclusion concerns a
single rational-prime component; no assertion about the sum over other
primes and the archimedean places follows from it.
\end{remark}

\section{Scope}\label{sec:scope}

\paragraph{Related work, added after the first deposit.}
\textup{[S]} \cite{LANA} isolates a compatibility problem in the passage from
\cite[Thm.~3.11]{IUT3} to \cite[Cor.~3.12]{IUT3}; its goal \cite[\S9.2]{LANA} is
$\eta^{q}=\eta^{\mathrm{anab}}_{S}$ for some suitable $S$, read in
\cite[\S9.3]{LANA} as the assertion that the rigidified $q$-pilot is represented
in the output regions, and it states that it has no proof of this at present.
\textup{[P]} That is not the statement \SA4, which concerns which regions the
union of possible images contains; within \cite{LANA} the corresponding place is
\cite[\S10.4]{LANA}.  \textup{[S]} That section says the hull computation ``is
made at each place individually, but the effects of $\Ind1$, $\Ind2$ and $\Ind3$
cannot be detected locally''.  \textup{[P]} This does not contradict the
conditional statement proved here; if that reading is correct then \SA{} fails and
the source application is withdrawn.  We have not identified the regions, hull and
readout of \cite{LANA} with the fixed component used here, and the note records
the cross-check as a falsifier.  One point of contact is exact: the LGP element of
\cite[\S5.2(f)]{LANA} carries $\underline{\underline q}_v^{\,j^2}$ in the
distinguished slot alone and $1$ in every other, the shape used here.
\cite{LANA} gives no explicit numerical initial $\Theta$-datum and computes no
fibre cardinality or constant/mixed classification of tuples, but its containers
are indexed as here \cite[\S5.2(a)]{LANA}, so the configuration is admitted by
that formalism.  \cite{Form} decomposes the passage as
$3.11\Rightarrow3.11.5\Rightarrow3.12$ by moving the ``hull+det'' of $3.12$ into
$3.11.5$ --- an explanatory label, which \cite{Form} states does not appear in
\cite{IUT1,IUT2,IUT3,IUT4} --- and we use it only for that.  We do not identify
the transport of \SA3 with the algorithmic parallel transport of
\cite[Rem.~3.11.1(iv)]{IUT3}, nor read that remark as a prohibition on a component
map between summands of one fixed ambient object.

\section{What would falsify the argument}\label{sec:falsify}

The interpretation-free content is Part~\ref{part:datum}: the datum of
Definition~\ref{def:datum} satisfies \cite[Definition 3.1]{IUT1}, with the local
profile of Proposition~\ref{prop:profile}.  That part stands or falls on exact
arithmetic and on the quoted conditions, and we invite direct checking.

The proposed application in the remainder rests on \SA{}.  Each of the
following would invalidate an identified step in our proposed derivation of
\SA{}, and require withdrawal of that application.  Such an objection need
not disprove the abstract conditional implication ``if \SA{}, then the
numerical comparison''; the issue is whether \SA{} holds for the source's
actual collection of possible images.
\begin{enumerate}[label=(F\arabic*)]
\item The transposition of two factors of a capsule with $\ge2$ factors is
  \emph{not} induced by an automorphism of the procession admitted in
  $\Ind1$ --- for instance if the distinguished slot is fixed by a condition we
  have not located.  This would contradict our reading of (SA2).
\item ``Regarded up to indeterminacies'' in \cite[Theorem 3.11(i)]{IUT3}, and
  ``subject to the indeterminacies'' in \cite[Corollary 3.12]{IUT3} and in the
  introduction of \cite{IUT3}, do \emph{not} mean that the $\Ind1$-orbit of a
  possible image consists of possible images.  This would contradict (SA4).
\item The specified native representation $\kappa_E[J]$ is not among the possible
  images of \cite[Corollary 3.12]{IUT3} --- for instance if a further condition
  on possible images, beyond interpretability as an
  $\mathcal{F}^\Vdash$-prime-strip, excludes it.  This would contradict (SA1).
\item The comparison of \cite[Theorem 3.11(ii)(a)]{IUT3} does not preserve the
  lattice $L$ at a tame place, contrary to the derivation in \S\ref{sec:sa}.
\end{enumerate}
We regard (F2) as the most likely locus of disagreement, since it is the point at
which the treatment of labels is at issue.  This is related to the label and
holomorphic-structure issue discussed in \cite[\S12]{Rpt2018}; that passage
does not itself decide the orbit-membership question posed here.

\section{What we ask of the reader}\label{sec:ask}

This paper is addressed, in the first instance, to those who know the theory
well.  We ask two specific things.

\medskip\noindent\textbf{First, check Theorem~\ref{thm:A}.}
It is a finite list of conditions about one explicit curve over a real quadratic
field.  It is independent of \SA{}, though it uses the cited standard reduction and
geometric results.  If a
condition of \cite[Definition 3.1]{IUT1} fails for this datum, the rest of the
paper is void and we would be grateful to be told which one.  The computations
and their reproducibility scope are listed in Appendix~\ref{app:comp}; the
geometric construction is checked separately in \S\ref{sec:def31}.

\medskip\noindent\textbf{Second, tell us which of (F1)--(F4) holds, if any.}
Theorem~\ref{thm:B} is not offered as a settled conclusion.  It is offered as the
consequence of a stated reading \SA{}, together with the request that the reading
be adjudicated.  \S\ref{sec:falsify} enumerates four ways its proposed derivation can fail, and
we regard (F2) --- whether the $\Ind1$-orbit of a possible image consists of
possible images --- as the likeliest.  A reply of the form ``(F2) holds, because
$\ldots$'' would settle the matter, and we would withdraw Theorem~\ref{thm:B}
accordingly.

\medskip\noindent
We have tried to make the paper easy to refute rather than hard to refute.  The
interpretive content is confined to one hypothesis, that hypothesis is stated in
one place, and the ways it can fail are enumerated.  If the estimate examined in
\S\ref{sec:local} can be repaired --- for instance by enlarging the intermediate
container and carrying the resulting loss of volume through the averaging of
\S\ref{sec:averaged} --- then that repair, and not this paper, is the interesting
mathematics, and we would like to see it.

\appendix

\section{The exact computations}\label{app:comp}

All computations of Part~\ref{part:datum} are exact integer computations in
$\mathbb{Z}$ or $\mathbb{Z}[\sqrt5]$.  We record the inputs so that they may be
repeated.
\begin{itemize}
\item $\pi=16+3\sqrt5$; $\Norm(\pi)=211$; $a=\pi^{29}=A+B\sqrt5$ with
  $A,B\in\mathbb{Z}$ of $40$ and $39$ decimal digits; $\Norm(a)=211^{29}$.
\item $s^2\equiv5\pmod{211}$ has solutions $s\equiv65,146$.
\item $211\nmid\Norm(1-a)$.  For reference,
  $\Norm(1-a)=2^2\cdot3^2\cdot5\cdot59\cdot80621\cdot C$ with
  $C$ a $59$-digit cofactor; the primality of $C$ is neither needed nor claimed.
\item The six orbit identities of Proposition~\ref{prop:fmod} all fail.
\item Traces of Frobenius $a_v$ for $121$ degree-one places of $F_0$ of good
  reduction, each satisfying the Hasse bound; the witness $(11,4,0)$ is used in
  Proposition~\ref{prop:sl2}.  The additional witnesses printed by the program
  are supplementary arithmetic checks.
\end{itemize}
The program \texttt{scripts/checks.py} checks these displayed computations
and the comparison constants.  The separate exhaustive integer sieve
\texttt{scripts/check\_norm\_seventh\_power\_free.cpp} verifies that
$\Norm(1-a)$ is seventh-power-free.  Its cofactor has exact seventh-root floor
$225469788$, and all $12407147$ primes up to that bound are tested; the listed
small factors are removed with their complete valuations first.  This proves
the assertion without a primality claim for the residual cofactor.
Neither program proves the cited geometric results or \SA{}.
The logarithm ratio in \S\ref{sec:averaged} is evaluated approximately for
orientation, while its required bound is proved exactly in the proof of
Theorem~\ref{thm:averaged}.
The companion \texttt{scripts/check\_crt\_family.py} reproduces the CRT
congruences and exact local orders for $N=1,106,211$.  These finite checks do
not certify any individual sample as seventh-power-free; the existence
assertion is proved by the counting argument in Appendix~\ref{app:family}.

\section{A machine-checked finite model}\label{app:lean}

The distinction of \S\ref{sec:twoparts} is elementary but easy to elide, so we
record it as a finite model checked by the Lean~4 kernel.  Subsets of a
four-element set are encoded as four-bit naturals; the action transposes the
first two and the last two points; \texttt{hull} fills the interval between the
least and greatest member; \texttt{card} counts members.  With
$A=\{0,2\}$ and $\sigma A=\{1,3\}$ the file establishes, with no axioms:
\begin{itemize}
\item \texttt{hull\_card\_invariant}: $\mathrm{card}(\mathrm{hull}\,A)
  =\mathrm{card}(\mathrm{hull}\,\sigma A)$ --- the analogue of (A);
\item \texttt{orbit\_not\_forced\_by\_invariance}: a predicate satisfying (A) but
  not (B) exists, so (A) does not imply (B);
\item \texttt{the\_point}: the orbit predicate satisfies both (A) and (B), and
  nevertheless $\mathrm{card}(\mathrm{hull}\,A)<
  \mathrm{card}\bigl(\mathrm{hull}(A\cup\sigma A)\bigr)$ and likewise for
  $\sigma A$.
\end{itemize}
The model is an analogy, not a formalization of the source; it fixes the logical
distinction only.

\section{Disclosure of the use of AI assistance}\label{app:ai}

Large language models were used substantially in the preparation of this paper:
in locating the relevant passages of \cite{IUT1,IUT3,IUT4,CanLift,MFHMP}, in
carrying out and cross-checking the computations, in drafting the exposition, and
in adversarial review of successive drafts.  Several distinct models were used,
and disagreements between them were resolved by returning to the sources.  This
note delimits what does and does not depend on that assistance.

\medskip\noindent\textbf{What does not depend on it.}
\begin{enumerate}[label=(\roman*)]
\item The displayed numerical calculations in Part~\ref{part:datum} are
  exact and reproducible by the programs described in Appendix~\ref{app:comp}.
  The remaining deductions use the supplied proofs and cited theorems about
  reduction, torsion and covers; they are not certified by the arithmetic
  programs.  None of these checks certifies the source reading \SA{}.
\item Every source condition invoked in \S\ref{sec:def31} and \S\ref{sec:sa} is
  cited by statement number and, where the wording matters, quoted verbatim.
\item The finite model of Appendix~\ref{app:lean} is checked by the Lean~4
  kernel and depends on no axioms.
\end{enumerate}

\medskip\noindent\textbf{What does depend on human judgement.}
\begin{enumerate}[label=(\roman*),start=4]
\item The hypothesis \SA{} of \S\ref{sec:sa} is a reading of the source.  It is
  not a computation and it is not a quotation; it is an interpretation of what
  ``regarded up to indeterminacies'' and ``subject to the indeterminacies''
  mean.  It is stated separately, its four supporting steps are given
  separately, and \S\ref{sec:falsify} states what would falsify it.  The author
  regards this as the point at which disagreement is most likely, and invites it.
\item The selection of which passages are the relevant ones, and the framing of
  the comparison in \S\ref{sec:local}, are likewise judgements.
\item An earlier draft stated Proposition~\ref{prop:exist} without a
  reconstructed proof.  The congruences and counting argument are now supplied
  in Appendix~\ref{app:family}, following source and arithmetic review with AI
  assistance.  This does not assert an independent human verification.  The
  proposition is used only in \S\ref{sec:averaged}, not in
  Theorem~\ref{thm:local}.
\end{enumerate}

\medskip\noindent
Errors found during preparation are recorded here rather than silently removed,
since the pattern is informative.  Three that were caught late: an appeal to a
Lean declaration absent from the checked source after a change of name; an
axiom tally whose itemized counts did not sum to its stated total; and a conjecture
that the four exceptional $j$-invariants of \S\ref{sec:kcore} would be algebraic
integers, which is false for two of them --- had it not been checked against
\cite[Proposition 2.1]{MFHMP}, Proposition~\ref{prop:kcore} would have reached a
correct conclusion by an incorrect route.  The author does not claim that no
further errors remain.  A subsequent audit of the accompanying short reports
corrected an omitted $j^2$ at stage $2$, an overly broad logarithm-lattice
identity, and a confusion between nonconstant place tuples and mixed reduction
types.  This paper already used $j^2N/7$ in its general-stage calculation;
version 0.1.2 makes the stage convention and small-ramification restriction
explicit throughout.  The labels [S], [P], [H] now distinguish source
statements, deductions or computations, and hypotheses, respectively.

\section{The seventh-power-free existence argument}\label{app:family}

\begin{proof}[Proof of Proposition~\ref{prop:exist}]
Write $p=211$.  Since $65^2\equiv5\pmod p$ and $130$ is a unit modulo $p$,
Hensel's lemma gives an integer $s_N$ satisfying
$s_N^2\equiv5\pmod{p^{N+1}}$ and $s_N\equiv65\pmod p$.
Choose $0\le A_0<M$ by the following complete set of CRT conditions:
\[
 \begin{split}
 A_0&\equiv s_N+p^N\pmod{p^{N+1}},\qquad A_0\equiv0\pmod8,\\
 A_0&\equiv2\pmod3,\quad A_0\equiv2\pmod5,\quad
 A_0\equiv2\pmod{19},\quad A_0\equiv6\pmod{11},\\
 M&=8\cdot3\cdot5\cdot11\cdot19\cdot p^{N+1}
   =25080p^{N+1}.
 \end{split}
\]
For $A=A_0+Mx$, $x\ge1$, put
\[
 f(T)=(T^2-5)((T-1)^2-5)
     =T^4-2T^3-9T^2+10T+20,\qquad
 G_N(x)=p^{-N}f(A_0+Mx).
\]
The constant term is divisible by $p^N$ and every positive-degree
coefficient by $M$, so $G_N\in\mathbb Z[x]$ has degree four.
The additional Hensel digit ensures exact order $N$, not merely order at
least $N$.  In fact, for every $x$,
\[
 \frac{A^2-5}{p^N}\equiv130\pmod p,\qquad
 (A-1)^2-5\equiv82\pmod p,\qquad G_N(x)\equiv110\pmod p.
\]
At the conjugate place $a\equiv130$, which is neither $0$ nor $1$.

The primes dividing $M$ cannot supply a seventh power:
$v_2(G_N(x))=2$, since $f(A)\equiv4\pmod8$;
for $q=3,5,19$, $f(A)\equiv f(2)=4\not\equiv0\pmod q$;
for $q=11$, $f(A)\equiv f(6)=620\equiv4\pmod{11}$;
and $p\nmid G_N(x)$ as just shown.
The two quadratic factors of $f$ have discriminant $20$ and resultant
$(1-2\sqrt5)(1+2\sqrt5)=-19$.  Hence
$\operatorname{disc}(f)=20^2\cdot19^2=144400$.
For any prime $q\nmid M$, all roots of $f$ modulo $q$ are simple, and there
are at most four.  Each lifts uniquely modulo $q^7$.
As $M$ and $p$ are units modulo $q$, at most four residue classes modulo
$q^7$ satisfy $q^7\mid G_N(x)$.

Fix $N$.  For $1\le x\le X$,
$0<G_N(x)<C_N(X+1)^4$, where one may take $C_N=M^4/p^N$.
Thus any prime whose seventh power divides $G_N(x)$ is at most
$Y=C_N^{1/7}(X+1)^{4/7}$.
The number of excluded integers $x$ is at most
\[
 \begin{split}
 \sum_{\substack{q\le Y\text{ prime}\\q\nmid M}}
       4\left(\frac{X}{q^7}+1\right)
 &\le 4X\sum_{n\ge7}\frac1{n^7}+4Y\\
 &\le\frac{X}{69984}+4C_N^{1/7}(X+1)^{4/7},
 \end{split}
\]
because $4\int_6^\infty t^{-7}\,dt=1/69984$.
The second term is $o(X)$ for this fixed $N$, while $1/69984<1$.
Therefore infinitely many $x$, and hence infinitely many $A$, remain.
The identity $f(A)=\Norm(a)\Norm(1-a)$ proves the assertion.
No bound uniform in $N$, and no certified individual value of $A$ for
$N=211$, is needed for this existence conclusion.
\end{proof}

The rational logarithm bounds used in Theorem~\ref{thm:averaged} are
proved there directly from the exponential series; they require no numerical
approximation to logarithms.

\begin{thebibliography}{MFHMP}

\bibitem[CanLift]{CanLift}
S.~Mochizuki,
\emph{The absolute anabelian geometry of canonical curves},
Doc. Math., Extra Volume: Kazuya Kato's Fiftieth Birthday (2003), 609--640.
\href{https://doi.org/10.4171/DMS/3/17}{doi:10.4171/DMS/3/17}.

\bibitem[EtTh]{EtTh}
S.~Mochizuki,
\emph{The \'etale theta function and its Frobenioid-theoretic manifestations},
Publ. Res. Inst. Math. Sci. \textbf{45} (2009), no.~1, 227--349.
\href{https://doi.org/10.2977/PRIMS/1234361159}{doi:10.2977/PRIMS/1234361159}.

\bibitem[GenEll]{GenEll}
S.~Mochizuki,
\emph{Arithmetic elliptic curves in general position},
Math. J. Okayama Univ. \textbf{52} (2010), 1--28.
\href{https://doi.org/10.18926/mjou/33503}{doi:10.18926/mjou/33503}.

\bibitem[IUT1]{IUT1}
S.~Mochizuki,
\emph{Inter-universal Teichm\"uller theory I: Construction of Hodge theaters},
RIMS Preprint 1756 (2012); consulted author version, May 2020.
Published in Publ. Res. Inst. Math. Sci. \textbf{57} (2021), no.~1/2, 3--207;
see also \cite{IUTpub}.
\href{https://doi.org/10.4171/PRIMS/57-1-1}{doi:10.4171/PRIMS/57-1-1}.

\bibitem[IUT2]{IUT2}
S.~Mochizuki,
\emph{Inter-universal Teichm\"uller theory II: Hodge--Arakelov-theoretic evaluation},
RIMS Preprint 1757 (2012); consulted author version, December 2020.
Published in Publ. Res. Inst. Math. Sci. \textbf{57} (2021), no.~1/2, 209--401;
see also \cite{IUTpub}.
\href{https://doi.org/10.4171/PRIMS/57-1-2}{doi:10.4171/PRIMS/57-1-2}.

\bibitem[IUT3]{IUT3}
S.~Mochizuki,
\emph{Inter-universal Teichm\"uller theory III: Canonical splittings of the log-theta-lattice},
RIMS Preprint 1758 (2012); consulted author version, May 2020.
Published in Publ. Res. Inst. Math. Sci. \textbf{57} (2021), no.~1/2, 403--626;
see also \cite{IUTpub}.
\href{https://doi.org/10.4171/PRIMS/57-1-3}{doi:10.4171/PRIMS/57-1-3}.

\bibitem[IUT4]{IUT4}
S.~Mochizuki,
\emph{Inter-universal Teichm\"uller theory IV: Log-volume computations and set-theoretic foundations},
RIMS Preprint 1759 (2012); consulted author version, April 2020.
Published in Publ. Res. Inst. Math. Sci. \textbf{57} (2021), no.~1/2, 627--723;
see also \cite{IUTpub}.
\href{https://doi.org/10.4171/PRIMS/57-1-4}{doi:10.4171/PRIMS/57-1-4}.

\bibitem[IUTpub]{IUTpub}
S.~Mochizuki,
\emph{Inter-universal Teichm\"uller theory I--IV},
collective reference to the four articles in
Publ. Res. Inst. Math. Sci. \textbf{57} (2021), no.~1/2:
I, 3--207; II, 209--401; III, 403--626; IV, 627--723.

\bibitem[Form]{Form}
S.~Mochizuki,
\emph{On the formalization of IUT: a preliminary progress report}
(joint work in progress with Y.~Hoshi, G.~Yamashita, Y.~Yang),
slides, Workshop on AI and Theorem Provers in Mathematics, April 2026, 14 pp.

\bibitem[LANA]{LANA}
Project LANA,
\emph{Project LANA interim report on IUT theory}, July 2026, 50 pp.
Lean for ANAbelian geometry; core members J.~Commelin, K.~Kedlaya, Y.~Hoshi,
A.~Topaz, F.~Kato (project leader).

\bibitem[MFHMP]{MFHMP}
S.~Mochizuki, I.~Fesenko, Y.~Hoshi, A.~Minamide and W.~Porowski,
\emph{Explicit estimates in inter-universal Teichm\"uller theory},
Kodai Math. J. \textbf{45} (2022), no.~2, 175--236.
\href{https://doi.org/10.2996/kmj45201}{doi:10.2996/kmj45201}.
Consulted author versions: December 2021 and May 2022.

\bibitem[Rpt2018]{Rpt2018}
S.~Mochizuki,
\emph{Report on discussions, held during the period March 15--20, 2018,
concerning inter-universal Teichm\"uller theory (IUTch)},
author manuscript, February 2019 revision, 45 pp.
\url{https://www.kurims.kyoto-u.ac.jp/~motizuki/Rpt2018.pdf}.

\bibitem[Sijs]{Sijs}
J.~Sijsling,
\emph{Canonical models of arithmetic $(1;\infty)$-curves},
in \emph{Arithmetic geometry: computation and applications},
Contemp. Math. \textbf{722}, Amer. Math. Soc., Providence, RI, 2019, 149--165.
\href{https://doi.org/10.1090/conm/722/14530}{doi:10.1090/conm/722/14530}.
Consulted author version: \href{https://arxiv.org/abs/1707.01158v2}{arXiv:1707.01158v2} (2017).

\bibitem[SS]{SS}
P.~Scholze and J.~Stix,
\emph{Why $abc$ is still a conjecture},
author manuscript, July 16, 2018, 10 pp.
\url{https://www.math.uni-bonn.de/people/scholze/WhyABCisStillaConjecture.pdf}.

\end{thebibliography}

\end{document}
